\chapter{Bloc 4 --- Data Lake OMEGA LAKE (C18-C21)}
\label{chap:bloc4}

Le Bloc 4 conçoit et opère le data lake OMEGA LAKE pour absorber la
donnée massive IoT de la flotte et des entrepôts. Clos le 23/09/2026
(milestone M3), les 4 compétences ont chacune un livrable vérifié --
détail exhaustif dans \texttt{docs/comptes\_rendus/M3.md}.

\section{C18 --- Concevoir l'architecture du \textit{data lake}}

La différence entre base de staging, entrepôt et data lake n'est pas le
volume ni le degré de nettoyage des données, mais le moment où le
schéma est décidé : \textit{schema-on-write} pour les deux premiers,
\textit{schema-on-read} pour le lake
(\texttt{docs/architecture/architecture\_omega\_lake.md}). Les flux IoT
transitent en flux séparé (jamais par la base de staging relationnelle),
mais la zone \texttt{curated/} est conçue pour être jointe à la volée
aux dimensions conformées de l'entrepôt via des clés naturelles réelles,
vérifiées dans les données plutôt que supposées.

\section{C19 --- Intégrer une infrastructure de données à partir de composants}

MinIO est déployé comme stockage objet compatible S3
(\texttt{docs/architecture/integration\_infrastructure\_omega\_lake.md}).
Le mécanisme de jointure --- DuckDB, choisi après vérification que le
projet n'utilise \texttt{pandas} nulle part --- est tranché explicitement
avant tout code de pipeline, et vérifié contre MinIO/PostgreSQL réels,
pas supposé équivalent à un test local :

\begin{lstlisting}[language=SQL,caption={Configuration DuckDB/\texttt{httpfs} pour MinIO (\texttt{src/datacore/storage/lake/transform.py})}]
SET s3_endpoint='{OMEGA_LAKE_S3_ENDPOINT}';
SET s3_access_key_id='{MINIO_ROOT_USER}';
SET s3_secret_access_key='{MINIO_ROOT_PASSWORD}';
SET s3_use_ssl=false;
SET s3_url_style='path';
\end{lstlisting}

L'ingestion de \texttt{raw/} passe par \texttt{boto3} (copie fidèle,
octet pour octet, vérifiée par somme de contrôle SHA-256) et non par
DuckDB, qui réinterpréterait le contenu. Le consommateur du flux SSE,
fenêtré par intervalle de 5~minutes plutôt que par jour, a été vérifié
par une coupure brutale réelle (\texttt{kill -9} en cours de réception
puis redémarrage) : aucune collision ni donnée corrompue.

\section{C20 --- Développer un catalogue de données}

Le périmètre de cette compétence a été explicitement élargi avant tout
code : le texte initial de l'issue ne parlait que de cataloguer « après
C19 », mais à ce stade les zones \texttt{staging/}/\texttt{curated/}
restaient vides -- construire un contenu réel y a été ajouté. Le
catalogue lui-même est généré par introspection réelle (\texttt{boto3}
+ \texttt{DESCRIBE} DuckDB), pas maintenu à la main, même principe que
l'audit RGPD automatisé de C16 : un catalogue généré ne peut pas
devenir silencieusement obsolète.

\section{C21 --- Garantir la gouvernance des données}

Le registre RGPD du lake (\texttt{docs/architecture/registre\_rgpd\_lake.md})
identifie un risque réel de ré-identification, trouvé en concevant et
non en auditant après coup : \texttt{vehicule\_id} (géolocalisation) se
recoupe, dans le temps, avec \texttt{tournees.chauffeur} de la base de
staging. Décision actée dès C18 : aucune jointure automatisée entre les
deux.

Une première analyse de pseudonymisation (« pas de protection réelle
possible, vu le faible nombre de véhicules ») s'est révélée incomplète
après relecture externe : vraie pour les Data Engineers (accès complet
des deux côtés), fausse pour le rôle \texttt{lake\_reader} (accès au
lake seul). Un pseudonyme HMAC à clé secrète --- un hash simple sans
clé aurait été recalculable par force brute sur les 15 véhicules du jeu
de données --- a été implémenté comme export dédié
(\texttt{curated\_bi/}), sans jamais modifier \texttt{curated/} ni
distribuer la clé au rôle de lecture :

\begin{lstlisting}[language=Python,caption={Garde-fou sur la clé de pseudonymisation (\texttt{src/datacore/storage/lake/curated\_bi.py})}]
def verifier_cle_configuree(cle: str | None) -> str:
    if not cle:
        raise RuntimeError("LAKE_PSEUDONYM_KEY n'est pas definie. ...")
    if cle in (VALEUR_EXEMPLE_ENV, ANCIENNE_VALEUR_REPLI_SUPPRIMEE):
        raise RuntimeError("LAKE_PSEUDONYM_KEY vaut une valeur publique connue. ...")
    if len(cle) < LONGUEUR_MINIMALE:
        raise RuntimeError(f"LAKE_PSEUDONYM_KEY est trop courte ({len(cle)} caracteres). ...")
    return cle
\end{lstlisting}

Ce garde-fou a lui-même été renforcé après qu'une faille a été trouvée
en seconde relecture externe --- un repli codé en dur pour cette clé
dans \texttt{config.py}, visible dans le dépôt public, aurait rendu la
pseudonymisation silencieusement inopérante --- puis à nouveau après
une question sur sa couverture au-delà du seul cas de la valeur
d'exemple. Les droits d'accès par groupe ont également été corrigés en
cours de route : la politique \texttt{lake\_reader} portait initialement
sur le bucket entier, permettant de contourner la pseudonymisation en
lisant \texttt{curated/} directement -- restreinte au seul préfixe
\texttt{curated\_bi/}, revérifiée bout en bout avec les vraies clés
\texttt{lake\_reader} (lecture de \texttt{curated\_bi/} réussit et
renvoie des données pseudonymisées ; la même requête contre
\texttt{curated/} échoue en HTTP~403). La procédure de purge RGPD est
vérifiée par un test réel : une partition synthétique délibérément
datée, injectée dans le lake réel, purgée avec succès et le résultat
confirmé indépendamment.
